%%%%%%%%%%%%%%%%%%%% manuscript.tex %%%%%%%%%%%%%%%%%%%%%%%%%%%%%%%%%
%
% ICIT/Springer svproc manuscript draft.
% This file intentionally preserves the svproc template structure while
% replacing the sample content with the shrimp disease segmentation paper.
%
%%%%%%%%%%%%%%%%%%%%%%%%%%%%%%%%%%%%%%%%%%%%%%%%%%%%%%%%%%%%%%%%%%%%%%

\documentclass{svproc}

% RECOMMENDED %%%%%%%%%%%%%%%%%%%%%%%%%%%%%%%%%%%%%%%%%%%%%%%%%%%%%%%%

\usepackage{url}
\usepackage{graphicx}
\def\UrlFont{\rmfamily}

\begin{document}
\mainmatter

\title{Specimen-Leakage-Aware Evaluation for Shrimp Disease Segmentation}

\titlerunning{Leakage-Aware Shrimp Disease Segmentation}

% TODO: Replace author placeholders before submission.
\author{Vinh Dinh Nguyen\inst{1} \and Phong Van Nguyen\inst{1} \and Nhan Huu Tran\inst{1} \and Nhu Huynh Le Thi\inst{1}}

\authorrunning{Vinh Dinh Nguyen et al.}

\tocauthor{Vinh Dinh Nguyen, Phong Van Nguyen, Nhan Huu Tran and Nhu Huynh Le Thi}

% TODO: Replace affiliation placeholders before submission.
\institute{Department of Artificial Intelligence, FPT University, Can Tho, Viet Nam\\
\email{vinhnd29@fe.edu.vn}
\email{phongnv.ce192081@gmail.com}
\email{nhanthce181655@fpt.edu.vn}
\email{nhuhuynhlethi1508@gmail.com}}


\maketitle

\begin{abstract}
Shrimp disease segmentation can support aquaculture monitoring by localizing visual disease regions rather than only assigning image-level labels. However, public shrimp disease images often contain multiple views of the same physical specimen, making random image-level train/test splits vulnerable to specimen leakage. This paper constructs an instance segmentation dataset from ShrimpDiseaseImageBD and evaluates YOLO-based instance segmentation under a specimen-grouped, disease-stratified protocol. We show that image-level splits create substantial train/test specimen overlap and higher test scores, while grouped splits provide a stricter estimate of unseen-specimen generalization. Under the leakage-free split, YOLO11n-seg gives the strongest accuracy and healthy-false-positive trade-off among the tested YOLO variants. We further report healthy-aware diagnostics and preliminary preprocessing/augmentation screening to guide future optimization.
\keywords{shrimp disease, instance segmentation, data leakage, grouped split, YOLO}
\end{abstract}

\section{Introduction}

Computer vision has become a practical tool for aquaculture health
monitoring, where early detection of visible disease symptoms can help
reduce economic loss and improve farm management. For shrimp disease,
image classification can indicate whether a disease is present, but
segmentation provides a more informative output by localizing the
visible diseased regions. This localization is important when symptoms
are small, spatially sparse, or mixed across disease types.

Reliable segmentation evaluation is difficult when the dataset contains
multiple images of the same physical shrimp specimen. If ordinary
image-level random splitting is used, one view of a shrimp may appear
in training while another view of the same shrimp appears in validation
or testing. A model can then benefit from specimen-specific appearance,
such as body texture, lighting, pose, and background, rather than
learning disease cues that generalize to unseen shrimp. This issue is
especially important in small biological datasets, where each specimen
contributes a nontrivial fraction of the available data.

This work addresses that problem for shrimp disease segmentation. We
derive a consistently hand-labeled segmentation dataset from the public
ShrimpDiseaseImageBD image dataset \cite{islam2025shrimpdiseaseimagebd},
define specimen identities from the filename convention, and enforce
a disease-stratified grouped split in which all images of the same
shrimp remain in exactly one partition. We then compare leakage-prone
image-level splits with leakage-free grouped splits, select a YOLO
segmentation backbone under the grouped protocol, and evaluate
preprocessing and augmentation strategies using both mAP and
healthy-aware diagnostics.
\section{Related Work}

\subsection{Shrimp Disease Image Datasets}

ShrimpDiseaseImageBD was released as a public dataset for
computer-vision-based shrimp disease detection in Bangladesh
\cite{islam2025shrimpdiseaseimagebd}. The dataset contains raw images
organized by health status and disease type, including healthy shrimp,
black gill (BG), white spot syndrome virus (WSSV), and co-infection
cases. Its public availability makes it a useful foundation for
classification and segmentation research, but its multi-image-per-specimen
structure requires careful splitting when evaluating generalization.
Recent shrimp-vision studies have explored transfer learning,
augmentation, explainability, and lightweight detection for shrimp
disease recognition \cite{jone2026shrimpxnet,fei2025lightweight}. These
works show the promise of image-based shrimp health monitoring, but they
primarily emphasize classification or detection accuracy. In contrast,
our study focuses on instance-level disease localization and on whether
the evaluation split is independent at the physical-specimen level.

\subsection{Instance Segmentation and YOLO}

YOLO-family models are widely used for real-time object detection and
segmentation because they provide a strong accuracy/speed trade-off and
are easy to train on custom datasets. In this work, we evaluate YOLOv8,
YOLO11, and YOLO26 segmentation variants using the Ultralytics training
framework \cite{ultralytics2026}. We focus on nano and small models
because shrimp disease segmentation is intended to remain practical for
iterative experimentation and future lightweight deployment.

\subsection{Data Leakage in Visual Learning}

Data leakage occurs when information unavailable at deployment time is
implicitly used during model development or evaluation
\cite{kaufman2012leakage}. In visual datasets, leakage can arise from
near-duplicate images, repeated subjects, or repeated specimens across
splits. The resulting test score can be overly optimistic because the
test images are no longer independent of the training images. Similar
concerns have been reported in medical imaging, where acquisition-site,
patient, or repeated-scan confounding can inflate apparent performance
\cite{zech2018confounding,roberts2021common,tampu2022inflation}. Our
setting is a specimen-level version of this problem: several images can
show the same shrimp, and all such images must be assigned to the same
split to estimate unseen-specimen performance.

\subsection{Preprocessing and Augmentation}

Small biological image datasets often benefit from augmentation and
preprocessing, but these methods must be evaluated under the same
leakage-free protocol as the baseline. We consider clean geometric and
photometric YOLO augmentations, copy-paste policies, Fourier high-pass
detail boosting, and wavelet-based enhancement as candidate methods.
Fourier preprocessing is motivated by the possibility that high-frequency
texture and boundary information may reveal subtle shrimp disease
patterns, but it is treated as an ablation rather than as an assumed
improvement.

\section{Dataset and Annotation}

\subsection{Source Dataset}

The image source is ShrimpDiseaseImageBD version 3
\cite{islam2025shrimpdiseaseimagebd}. The Mendeley record reports 1149
raw images across four disease-status folders: 403 healthy, 198 BG, 328
WSSV, and 220 BG\_WSSV images. These images were captured from 416
shrimp specimens, ranging from one to three images per specimen. The
original filename convention has the
form
\begin{center}
\texttt{<Disease>-<ShrimpID>-img-<ImageNumber>.jpg}.
\end{center}
We use this structure to infer specimen groups as
\texttt{<Disease>::<ShrimpID>} and use this identifier as the specimen
group.

\subsection{Hand-Labeled Segmentation Dataset}

The segmentation dataset used in this paper was relabeled by a single
annotator to reduce inter-annotator variation in mask placement,
granularity, and disease-region interpretation. This decision is
important because disease segmentation masks are more sensitive to
annotation style than image-level labels. Inconsistent boundaries,
fragmented masks, and incomplete disease labels can directly change
mask mAP and healthy false-positive behavior. For repeated images of the same specimen, disease regions were annotated according to anatomical location rather than image position. 

Figure~\ref{fig:hand-labeled} visualizes how images of the same specimen differ and 
are segmented consistently across changes in position, pose, and capture
angle. 
\begin{figure}
    \centering
    \includegraphics[width=0.8\linewidth]{hand-labeled.png}
    \caption{Consistent labels for multiple images of the same specimen.}
    \label{fig:hand-labeled}
\end{figure}

Table~\ref{tab:dataset_stats} summarizes the analyzed hand-labeled
dataset. Empty label files represent healthy images, while diseased
images can contain BG masks, WSSV masks, or both.

\begin{table}
\caption{Dataset statistics for the hand-labeled shrimp segmentation set.}
\label{tab:dataset_stats}
\begin{center}
\begin{tabular}{lr}
\hline
Property & Count \\
\hline
Images & 1149 \\
Specimens & 416 \\
Labeled diseased images & 746 \\
Healthy or empty-label images & 403 \\
Mask instances & 1031 \\
BG mask instances & 462 \\
WSSV mask instances & 569 \\
Co-infection images & 220 \\
Maximum images per specimen & 3 \\
\hline
\end{tabular}
\end{center}
\end{table}

% TODO: Add a figure with four panels: healthy, BG, WSSV, and BG+WSSV
% overlays. Suggested label:
% \begin{figure}
% \begin{center}
% \includegraphics[width=0.95\textwidth]{figures/dataset_examples.pdf}
% \caption{Example hand-labeled shrimp disease segmentation masks.}
% \label{fig:dataset_examples}
% \end{center}
% \end{figure}

\section{Specimen-Grouped Evaluation Protocol}

\subsection{Problem Definition}

Let $G_i$ denote the specimen group for image $i$. A leakage-free split
must satisfy
\begin{equation}
G_i = G_j \Rightarrow S_i = S_j,
\end{equation}
where $S_i$ is the assigned split for image $i$. In other words, all
images of the same shrimp specimen must be placed together in training,
validation, or testing. If this constraint is violated, the test set may
contain shrimp that the model has already seen from another view during
training.

\subsection{Split Policies}

We compare three split policies:
\begin{itemize}
\item \textbf{Plain random image split}: images are shuffled and split
without disease stratification or specimen grouping.
\item \textbf{Stratified random image split}: image-level splitting
preserves disease distribution but still allows specimen leakage.
\item \textbf{Stratified grouped-specimen split}: specimens are grouped
first, stratified by disease status, and assigned so that train/test
specimen overlap is zero.
\end{itemize}

The final paper-facing protocol is the stratified grouped-specimen split.
The random image splits are retained only as leakage-prone references
for quantifying how much image-level evaluation can overstate
performance. The stratified random image split separates the effect of
class-distribution preservation from the effect of specimen grouping.

\subsection{Leakage Verification}

For every run, we save a split manifest, split summary, and split
fingerprint. We also compute train/test and validation/test specimen
overlap. A grouped split is accepted only when both overlap values are
zero. This makes the leakage constraint directly auditable rather than
being an informal assumption.

\section{Experimental Setup}

\subsection{Training Configuration}

All main experiments use instance segmentation models from the YOLO
family. Models are trained at image size 640 for
up to 100 epochs with patience 30. The default paper baseline uses the
Ultralytics augmentation hook and a clean-light YOLO augmentation policy:
horizontal flip probability 0.5, HSV perturbation
$(h=0.01,s=0.35,v=0.20)$, translation 0.05, scale 0.20, and no mosaic,
mixup, cutmix, or copy-paste. The best checkpoint is selected using
validation metrics only, and all test metrics are reported after
checkpoint selection. We use this simple default training configuration
to establish a fair baseline, select one backbone for later optimization,
and keep subsequent preprocessing or augmentation gains interpretable.

\subsection{Evaluation Sets}

We report metrics on two test views:
\begin{itemize}
\item \textbf{Full test}: includes both diseased and healthy images.
\item \textbf{Labeled-only test}: includes only images with disease
masks, allowing disease localization quality to be measured without
healthy background images diluting the instance metrics.
\end{itemize}
This separation is necessary because practical disease monitoring
requires both strong disease localization and low false-positive
prediction on healthy shrimp.

\subsection{Metrics}

The primary segmentation metrics are mask mAP50 and mask mAP50-95. We
also record box mAP for completeness. To evaluate practical behavior on
healthy shrimp and multi-mask disease images, we additionally report:
\begin{itemize}
\item healthy mask false-positive rate, the fraction of healthy images
with at least one predicted disease mask;
\item healthy false-positive masks per image;
\item disease mask missed-image rate, the fraction of diseased images
where no disease mask is predicted;
\item mask count mean absolute error (MAE), measuring disagreement
between predicted and ground-truth mask counts.
\end{itemize}

For validation-based model selection, we use a healthy-aware score:
\begin{equation}
S = \mathrm{mAP50}_{lab}
- 0.05 \cdot \mathrm{MAE}_{count}
- 0.15 \cdot R_{miss}
- 0.10 \cdot R_{healthyFP},
\end{equation}
where $\mathrm{mAP50}_{lab}$ is labeled-only mask mAP50,
$R_{miss}$ is the disease missed-image rate, and $R_{healthyFP}$ is the
healthy false-positive rate. Test-set healthy-aware scores are reported
only for analysis and are not used for model selection. The weights keep
mask mAP50 as the dominant term while penalizing clinically undesirable
missed diseased images, poor mask-count behavior, and false alarms on
healthy shrimp.

\section{Results}

\subsection{Three-Seed Split Policy Comparison}
To test whether the split-policy conclusion is stable, we repeat the
comparison using seeds 42, 123, and 3407.

Table~\ref{tab:split_policy_summary} summarizes YOLO11n-seg across
three seeds (42, 123, and 3407) under the three split policies. Both
image-level policies retain substantial train/test specimen overlap,
with 98.7 overlapping specimen groups for plain random splitting and
94.0 for stratified random splitting on average. In contrast, the
stratified grouped-specimen split keeps train/test overlap at zero for
every seed. The image-level splits obtain higher mAP values, but these
scores are not independent estimates of unseen-specimen performance
because test specimens are often represented by related images in
training.

\begin{table}
\caption{Three-seed split-policy summary using YOLO11n-seg. Values are mean $\pm$ standard deviation.}
\label{tab:split_policy_summary}
\begin{center}
\scriptsize
\begin{tabular}{lrrrrrr}
\hline
Split policy & Overlap & Full mAP50 & Labeled mAP50 & Labeled mAP50-95 & Healthy FP & Healthy-aware \\
\hline
Plain random image & $98.7 \pm 1.2$ & $0.671 \pm 0.013$ & $0.677 \pm 0.013$ & $0.299 \pm 0.024$ & $0.054 \pm 0.027$ & $0.648 \pm 0.012$ \\
Stratified random image & $94.0 \pm 6.1$ & $0.651 \pm 0.060$ & $0.659 \pm 0.062$ & $0.282 \pm 0.016$ & $0.098 \pm 0.024$ & $0.628 \pm 0.053$ \\
Stratified grouped specimen & $0.0 \pm 0.0$ & $0.451 \pm 0.027$ & $0.473 \pm 0.035$ & $0.144 \pm 0.021$ & $0.236 \pm 0.110$ & $0.415 \pm 0.039$ \\
\hline
\end{tabular}
\end{center}
\end{table}
Stratified random splitting has a lower mean score than plain random
splitting and higher variance across seeds. This suggests that
distribution balancing alone does not remove specimen leakage and that
image-level split performance remains sensitive to which related
specimen views are assigned to each partition.
The grouped score is lower, but it is more faithful to the deployment
question: how well does the model generalize to shrimp specimens that
were not observed during training? The grouped protocol also shows that
small leakage-free datasets still have meaningful split variance, so
final claims should report seed summaries rather than relying on a
single split. 

\begin{table}
\caption{Per-seed split-policy results using YOLO11n-seg.}
\label{tab:split_policy_per_seed}
\begin{center}
\scriptsize
\begin{tabular}{llrrrrr}
\hline
Split policy & Seed & Overlap & Full mAP50 & Labeled mAP50 & Healthy FP & Healthy-aware \\
\hline
Plain random image & 42 & 98 & 0.658 & 0.664 & 0.068 & 0.636 \\
Plain random image & 123 & 100 & 0.685 & 0.690 & 0.071 & 0.660 \\
Plain random image & 3407 & 98 & 0.671 & 0.676 & 0.023 & 0.647 \\
Stratified random image & 42 & 98 & 0.584 & 0.589 & 0.073 & 0.567 \\
Stratified random image & 123 & 87 & 0.699 & 0.706 & 0.098 & 0.666 \\
Stratified random image & 3407 & 97 & 0.669 & 0.682 & 0.122 & 0.650 \\
Stratified grouped specimen & 42 & 0 & 0.482 & 0.512 & 0.244 & 0.460 \\
Stratified grouped specimen & 123 & 0 & 0.440 & 0.462 & 0.341 & 0.388 \\
Stratified grouped specimen & 3407 & 0 & 0.432 & 0.445 & 0.122 & 0.399 \\
\hline
\end{tabular}
\end{center}
\end{table}

\subsection{YOLO Model Selection Under the Fair Split}
Before studying preprocessing and augmentation, we first choose a fixed
segmentation backbone under the leakage-free grouped protocol. We train
the same baseline configuration across YOLOv8-seg, YOLO11-seg, and
YOLO26-seg nano/small variants.
Table~\ref{tab:model_sweep} reports the grouped-specimen model sweep.
YOLO11n-seg provides the best labeled-only test mask mAP50, the best
full-test mask mAP50, and the best healthy-aware test score among the
tested variants. YOLO26s-seg gives the highest mAP50-95, and YOLO26n-seg
has the lowest healthy false-positive rate, but neither provides the
best overall trade-off. We therefore use YOLO11n-seg as the fixed
backbone for subsequent preprocessing and augmentation experiments.

\begin{table}
\caption{YOLO model sweep under the stratified grouped-specimen split.}
\label{tab:model_sweep}
\begin{center}
\scriptsize
\begin{tabular}{lrrrrrr}
\hline
Model & Params (M) & Full mAP50 & Labeled mAP50 & mAP50-95 & Healthy FP & Healthy-aware \\
\hline
YOLO11n-seg & 2.877 & 0.482 & 0.512 & 0.168 & 0.244 & 0.460 \\
YOLO26s-seg & 11.506 & 0.417 & 0.460 & 0.182 & 0.268 & 0.397 \\
YOLOv8n-seg & 3.410 & 0.395 & 0.432 & 0.139 & 0.341 & 0.353 \\
YOLOv8s-seg & 11.821 & 0.391 & 0.444 & 0.159 & 0.341 & 0.364 \\
YOLO26n-seg & 3.126 & 0.430 & 0.463 & 0.160 & 0.195 & 0.393 \\
YOLO11s-seg & 10.113 & 0.297 & 0.361 & 0.124 & 0.512 & 0.250 \\
\hline
\end{tabular}
\end{center}
\end{table}

\subsection{Preprocessing and Augmentation Screening}

Preprocessing and augmentation experiments screen several lightweight
ways to improve the grouped YOLO11n-seg baseline: copy-paste, train-only
photometric augmentation, Fourier high-pass preprocessing, and
wavelet-based enhancement. Fourier and wavelet variants are evaluated
either as train-only transformed copies or as deterministic preprocessing
applied to all splits. Table~\ref{tab:preproc_partial} summarizes the
representative completed rows. None improves the healthy-aware baseline:
some methods slightly increase labeled-only mAP50, but the gain is
offset by lower mAP50-95, lower full-test mAP50, or higher false
positives on healthy shrimp.

\begin{table}
\caption{Preliminary preprocessing and augmentation screening.}
\label{tab:preproc_partial}
\begin{center}
\scriptsize
\begin{tabular}{lrrrrr}
\hline
Method & Full mAP50 & Labeled mAP50 & mAP50-95 & Healthy FP & Healthy-aware \\
\hline
YOLO11n clean-light baseline & 0.482 & 0.512 & 0.168 & 0.244 & 0.460 \\
Copy-paste 0.10 mixup & 0.464 & 0.515 & 0.144 & 0.488 & 0.424 \\
Photometric medium, train only & 0.443 & 0.477 & 0.168 & 0.244 & 0.420 \\
Noise/blur/compression, train only & 0.449 & 0.470 & 0.136 & 0.366 & 0.401 \\
Fourier s50 a0.30, hook on, train only & 0.436 & 0.468 & 0.163 & 0.341 & 0.403 \\
Fourier s50 a0.30, hook off, all splits & 0.452 & 0.498 & 0.166 & 0.244 & 0.426 \\
Wavelet db2, hook on, all splits & 0.406 & 0.435 & 0.131 & 0.293 & 0.352 \\
\hline
\end{tabular}
\end{center}
\end{table}

\section{Discussion}

The experiments support three observations. First, specimen leakage is
not a minor implementation detail. In the three-seed split-policy study,
plain random and stratified random image-level splits place 98.7 and
94.0 specimen groups across train and test on average, whereas the
grouped split eliminates this overlap for every seed. The corresponding
metric drop should be interpreted as a correction toward fairer
unseen-specimen evaluation, not as a failure of the model.

Second, stratified grouped splitting has an unavoidable cost on small
datasets. Once all images of a shrimp must move together, the number of
possible balanced train/validation/test combinations decreases. This
can make the test set harder or easier depending on which specimens are
assigned to it. The three-seed results show that leakage-free
evaluation is fairer, but still needs seed reporting or confidence
intervals when dataset size is limited.

Third, a strong default YOLO baseline is difficult to improve with
simple preprocessing. The default augmentation hook already includes
low-probability blur, grayscale, and CLAHE-style operations. Additional
photometric transforms or frequency-domain enhancement may help under
specific settings, but they can also increase healthy false positives or
overfit validation appearance. This justifies the use of healthy-aware
metrics rather than selecting only by labeled-test mAP.

The code for this study will be released at
\url{https://github.com/your-org/shrimp-leakage-aware-segmentation}.

The labeled instance segmentation dataset is publicly available as a Roboflow Universe dataset \cite{nguyen2026shrimpseg}; Roboflow was used for annotation management and YOLO-format export \cite{roboflow}.
\section{Conclusion}

This paper presents a leakage-aware evaluation study for shrimp disease
instance segmentation. We construct a single-annotator segmentation
dataset from ShrimpDiseaseImageBD, define specimen identities from the
filename convention, and enforce a disease-stratified grouped split to
prevent same-shrimp train/test leakage. The split-policy experiments
show that image-level splitting creates substantial specimen overlap and
higher apparent test performance. Under the leakage-free protocol,
YOLO11n-seg provides the strongest baseline among the tested YOLO
variants. Ongoing preprocessing and augmentation experiments will be
reported as secondary optimization results, with selection guided by
both segmentation mAP and healthy false-positive behavior.

% ---- Bibliography ----

\bibliographystyle{splncs03_unsrt}
\bibliography{references}

\end{document}
